\section{The total-spin limit}\label{sec:global-spin}

We now combine spatial comparison with the identification of equilibrium
magnetization lengths. Fix a finite $\beta$ above all the thresholds in
the preceding sections and write $m=m_{S,\beta}>0$.

\begin{lemma}\label{lem:local-second-moment}
For every finite nonempty rectangle $C\subset\Z^3$ and every
translation-invariant equilibrium state $\omega$,
\begin{equation}\label{eq:local-second-moment}
 \sum_{a=1}^3\omega\!\left[
       \left(\frac1{|C|}\sum_{x\in C}S_x^a\right)^2\right]\ge m^2.
\end{equation}
\end{lemma}
\begin{proof}
In a translation-ergodic equilibrium state, translation invariance and
positivity of variance give
\[
 \omega\!\left[
       \left(\frac1{|C|}\sum_{x\in C}S_x^a\right)^2\right]
    \ge\omega(S_0^a)^2.
\]
Summing and applying Proposition~\ref{prop:equilibrium-radius} proves
the assertion for extreme equilibrium states. The equilibrium set is
a compact face of the translation-invariant state space, as recalled in
Section~\ref{sec:equilibrium}. By the Krein--Milman theorem it is the
closed convex hull of its extreme points. The left side of
\eqref{eq:local-second-moment} is continuous and affine in the state,
so the inequality holds throughout that face.
\end{proof}

\begin{proposition}\label{prop:global-second-moment}
The periodic zero-field Gibbs states satisfy
\begin{equation}\label{eq:global-second-lower}
 \liminf_{L\to\infty}\frac1{V^2}
        \sum_{a=1}^3\langle(M_L^a)^2\rangle\ge m^2.
\end{equation}
\end{proposition}
\begin{proof}
Choose a sequence of even side lengths realizing the lower limit in
\eqref{eq:global-second-lower}, and pass to a subsequence on which the
Gibbs states converge locally. Their limit
$\omega$ is a translation-invariant equilibrium state by
Section~\ref{sec:equilibrium}. Fix the terminal scale $R_*$ in
Proposition~\ref{prop:scale-comparison}. Its finest partition consists
of rectangles with side lengths bounded above and below by constants
depending on $R_*$ only. For this fixed $R_*$ there are finitely many
possible rectangle shapes. Translation invariance and local convergence
therefore imply, uniformly over all cells in that partition,
\[
 \liminf_{L\to\infty}
 \frac1V\sum_{C\in\mathcal P_*}|C|
       \sum_{a=1}^3
       \left\langle
          \left(\frac1{|C|}\sum_{x\in C}S_x^a\right)^2
       \right\rangle
       \ge m^2,
\]
where Lemma~\ref{lem:local-second-moment} supplies the lower bound for
each limiting shape. The spatial comparison
\eqref{eq:spatial-comparison}, applied in each of the three directions,
gives
\[
 \liminf_{L\to\infty}\frac1{V^2}
       \sum_{a=1}^3\langle(M_L^a)^2\rangle
       \ge m^2-\frac{C_\beta}{\log R_*}.
\]
Let $R_*\to\infty$. The sequence was chosen to realize the lower limit,
so \eqref{eq:global-second-lower} follows along all even side lengths.
\end{proof}

\begin{proof}[Proof of Theorem~\ref{thm:main}]
The Hamiltonian commutes with the global spin representation of
$SU(2)$. Decompose the finite-volume Hilbert space into spin-$J$
irreducible representations and their multiplicity spaces. On each
irreducible factor the Gibbs density is a scalar multiple of the
identity. Thus the Gibbs trace defines a random label $J_L$ with
$0\le J_L\le SV$, and
\begin{equation}\label{eq:casimir}
 \sum_{a=1}^3\langle(M_L^a)^2\rangle
       =\E[J_L(J_L+1)].
\end{equation}
Conditional on $J_L=J$, a component in any fixed unit direction has the
uniform distribution on $-J,-J+1,\ldots,J$.

We first verify the upper-tail claim \eqref{eq:overview-upper}. Fix
$\epsilon>0$. By the definition of the right derivative, some fixed
$h>0$ satisfies
\[
 p_{S,\beta}(h)-p_{S,\beta}(0)
       \le h(m+\epsilon/4).
\]
Since $H_{S,L}$ commutes with $M_L^3$, exponential Markov and the pressure
limit give
\begin{align*}
 \limsup_{L\to\infty}\frac1V
   \log\PP\bigl(M_L^3/V>m+\epsilon/2\bigr)
 &\le -\beta h(m+\epsilon/2)
       +\beta\bigl[p_{S,\beta}(h)-p_{S,\beta}(0)\bigr] \\
 &\le -\beta h\epsilon/4<0.
\end{align*}
Here the probability is the spectral distribution of $M_L^3$ in the
Gibbs state. If $J/V>m+\epsilon$, at least $\epsilon V/2-1$ of its
$2J+1$ magnetic levels exceed $(m+\epsilon/2)V$. For all sufficiently
large $L$, their conditional probability is bounded below by a positive
constant depending only on $\epsilon,S$. Consequently
\[
 \PP(J_L/V>m+\epsilon)\longrightarrow0.
\]

The laws of $J_L/V$ are supported in the compact interval $[0,S]$.
Every subsequential limit is supported in $[0,m]$ by this upper-tail
bound. By \eqref{eq:casimir}, Proposition~\ref{prop:global-second-moment},
and $J_L/V^2\le S/V$, its second moment is at least $m^2$. A probability
measure on $[0,m]$ with that property is the point mass at $m$. Hence
\begin{equation}\label{eq:spin-concentration}
 J_L/V\longrightarrow m\quad\text{in probability},
 \qquad
 \frac{\E[J_L(J_L+1)]}{V^2}\longrightarrow m^2.
\end{equation}
The second conclusion proves \eqref{eq:second-moment-law}.

For $t\ne0$, rotate the direction $t/|t|$ to the third axis. Conditional
on $J_L=J$, the Laplace transform in \eqref{eq:spherical-law} is
\[
 \frac1{2J+1}\sum_{r=-J}^{J}e^{|t|r/V}.
\]
Whenever $J/V\to m>0$, these Riemann sums converge to
\[
 \frac12\int_{-1}^1e^{m|t|u}\dd u
       =\frac{\sinh(m|t|)}{m|t|}.
\]
They are uniformly bounded by $e^{S|t|}$, so
\eqref{eq:spin-concentration} implies convergence of their expectations.
The case $t=0$ is immediate. The normalized area measure on
$\mathbb S^2$ has a uniform third coordinate in $[-1,1]$, which identifies
the integral in \eqref{eq:spherical-law} and completes the proof.
\end{proof}
